% Internal Chinese methods draft. Official LIPIcs class is kept unmodified.
% Compile with ./build.sh; Chinese support requires XeLaTeX.
\documentclass[a4paper,UKenglish,cleveref,autoref,thm-restate]{lipics-v2021}

\usepackage[UTF8,scheme=plain,fontset=fandol]{ctex}
\usepackage{booktabs}
\usepackage{algorithm}
\usepackage{algpseudocode}
\usepackage{placeins}
\graphicspath{{template/}}
\bibliographystyle{plainurl}

\floatname{algorithm}{算法}
\algrenewcommand{\algorithmicrequire}{\textbf{输入：}}
\algrenewcommand{\algorithmicensure}{\textbf{输出：}}
\renewcommand{\proofname}{证明}

\title{可行见证驱动的非线性整数优化}
\subtitle{约束耦合搜索、核心反馈与整数射线认证}
\titlerunning{可行见证驱动的非线性整数优化}
% No author identities, affiliations, funding, or acceptance metadata are assumed.
\author{方法工作稿}{内部研究讨论，非投稿版本}{}{}{}
\authorrunning{OMT(NIA) 方法工作稿}
\Copyright{作者信息尚未设置}
\ccsdesc[500]{Theory of computation~Logic and verification}
\keywords{非线性整数优化, 随机局部搜索, 冲突核心, 无界证书, SMT}
\category{中文方法工作稿，2026-11-02}
\hideLIPIcs
\nolinenumbers

\begin{document}
\maketitle

\begin{abstract}
非线性整数优化既需要在离散可行域中持续改善目标，也需要区分当前可行值、有限最优值与无界情形。本文围绕可行见证组织搜索与认证：首先，以约束依赖和因子覆盖构造能够联合恢复可行性的目标改善动作；其次，利用锚点不可满足核心扩大改善邻域，并向局部搜索反馈变量优先级；最后，通过原式严格改善检查确认有限最优，并从可行见证合成可精确核验的整数改善射线以认证无界。整体框架只接纳经过完整精确检查的改善候选，并将局部失败与全局认证明确区分。可行见证由此同时服务于邻域构造与证书合成，将有限改善、有限最优和无界三类证据连接在同一优化框架中。
\end{abstract}

\section{整体框架}
\label{sec:overview}

我们考虑具有整数多项式约束的单目标优化问题。为统一记号，以下以最小化为例；最大化可对目标取负后对偶处理。输入问题记为 $(F_s,f_s)$，目标保持的预处理将其转换为搜索表示 $(F,f)$，其中 $F=\bigwedge_j C_j$ 为整数算术文字和布尔文字组成的 CNF。当前最好可行赋值记为 $a$，目标值为 $b=f(a)$。严格改善问题分别记为
\[
 Q_b=F\land f<b,\qquad Q_b^s=F_s\land f_s<b.
\]
算法~\ref{alg:overall} 给出整体流程。前端若已确定硬约束不可满足，则直接报告不可行；伪代码展开其余情况下的搜索与认证过程。

\begin{algorithm}[htbp]
\caption{可行见证驱动的搜索与认证}
\label{alg:overall}
\begin{algorithmic}[1]
\Require 原约束 $F_s$，整数目标 $f_s$，总截止时间 $D$
\Ensure 有限最优及见证、无界证书、最好可行见证，或未找到见证
\State $(F,f)\gets\Call{TargetPreservingImport}{F_s,f_s}$
\State $a\gets\bot$; $b\gets+\infty$; 初始化搜索状态 $\sigma$、邻域状态 $\eta$、证明状态 $\Pi$、优先级 $H$
\While{$\mathrm{now}<D$}
  \If{$a=\bot$}
    \State $e\gets\Call{FeasibilityTurn}{F_s,f_s,F,f,\sigma,D}$
  \Else
    \State $(k,\tau)\gets\Call{Schedule}{b,\text{历史},D}$
    \If{$k=\mathsf{Search}$}
      \State $e\gets\Call{CoupledImprove}{F,f,a,b,\sigma,H,\tau}$ \Comment{C1，\S\ref{sec:coupled}}
    \ElsIf{$k=\mathsf{Core}$}
      \State $e\gets\Call{CoreNeighborhood}{F,f,a,b,\eta,\tau}$ \Comment{C2，\S\ref{sec:core}}
    \ElsIf{$k=\mathsf{Ray}$}
      \State $e\gets\Call{WitnessRay}{F,f,a,b,\tau}$ \Comment{C3，\S\ref{sec:certificate}}
    \Else
      \State $e\gets\Call{SourceStrictCheck}{F_s,f_s,b,a,\Pi,\tau}$
    \EndIf
  \EndIf
  \If{$\mathrm{now}\ge D$} \State \textbf{break} \EndIf
  \State $D_e\gets e$ 的适用截止（缺省为 $D$）; $H\gets H\cup e.\mathit{corevars}$
  \If{$e$ 含当前界 $b$ 的原式严格改善查询 SAT 证据}
    \State 标记此界已被反驳，$b$ 改变前暂停同界 UNSAT 尝试
  \EndIf
  \If{$e$ 含候选 $c$}
    \If{$\Call{ExactAccept}{F,f,c,b,\min(D,D_e)}$}
      \State $(a,b)\gets(c,f(c))$; 更新起点，清理旧界状态，$H\gets\varnothing$
    \EndIf
  \ElsIf{$e=\mathsf{OptimalAt}(b)$ 且 $\mathrm{now}<\min(D,D_e)$}
    \State \Return $\mathsf{Optimal}(a,b)$
  \ElsIf{$e=\mathsf{VerifiedRay}(R)$ 且 $\mathrm{now}<\min(D,D_e)$}
    \State \Return $\mathsf{Unbounded}(R)$
  \EndIf
  \State 更新调度历史，保留可续传状态
\EndWhile
\State \Return 若 $a\ne\bot$ 则 $\mathsf{Feasible}(a,b)$，否则 $\mathsf{NoWitness}$
\end{algorithmic}
\end{algorithm}

整体框架以一个独立保存、经过精确检查的可行见证连接三个机制。\emph{约束耦合的目标改善搜索（C1）}将目标探测与关联变量修复组成联合动作；\emph{冲突核心驱动的邻域扩展与反馈（C2）}识别妨碍改善的固定条件，并将相关变量反馈给搜索；\emph{可行见证驱动的认证（C3）}在原式严格改善检查之外，从一个可行点合成永久改善的整数射线。三个机制及原式证明通道按条件交替运行，并非在每一轮依次全部执行。

\subsection{探索状态、可行见证与认证状态}
探索状态 $\sigma$ 允许暂时违反约束，$(a,b)$ 则始终保存已经校验的最好可行见证。$\eta$ 和 $\Pi$ 分别维护特定目标界下的邻域与证明进度。\textsc{ExactAccept} 从候选独立重算完整 $F$ 与 $f$，只接纳可行且严格改善的赋值；因此，已报告目标值单调下降。一次改善使旧界相关状态和核心优先级失效，但无需丢弃全部搜索信息。

初始阶段结合渐进式初始化、可恢复的 NIA 局部搜索和原式可行性救援。获得初解后，搜索阶段负责产生更好的候选，认证阶段负责区分当前值只是可行值，还是已经有限最优或无界。所有阶段共享截止时间 $D$。若结果带有更早的接纳截止 $D_e$，其重验与发布也必须在 $\min(D,D_e)$ 前完成；子过程时间片耗尽只表示暂未得到可用结论。

\subsection{两个表示之间的语义前提}
后续可靠性论证以前端保持可实现目标值集合为前提：
\begin{equation}
 \{f_s(m)\mid F_s(m)\}=\{f(x)\mid F(x)\}.
 \label{eq:values}
\end{equation}
搜索见证 $a$ 属于预处理后的表示，不能直接视为已经恢复全部原式变量的模型。普通候选和射线的精确检查针对 $(F,f)$；原式 SAT 模型则通过辅助变量补全过程映射到搜索表示。式~\eqref{eq:values} 表达导入应保持的语义关系，并非另有一个运行时等价性证明器。

\FloatBarrier
\section{约束耦合的目标改善搜索}
\label{sec:coupled}

\subsection{将目标探测与可行性恢复结合}
非线性整数约束将多个变量的改变耦合在一起。仅改变一个目标相关坐标，可能立即破坏等式或乘积关系，使一个有希望的改善方向被拒绝。我们将目标探测和关联变量修复组织为联合动作：先提出希望保留的目标变化，再寻找能够恢复硬约束的配套坐标变化。

搜索维护探索赋值 $x$，而非强制每一步停留在当前最好赋值 $a$。候选先按加权未满足子句代价比较，再按修复距离和目标值比较。在可行区域内，普通动作要求保持硬约束并改善目标；离开可行域后，搜索以恢复约束为主，同时保留已经取得的严格改善方向。对于尚未优于 $b$ 的状态，虚拟严格改善惩罚使搜索避免反复返回已经取得的目标水平。

与这一评分相配合，私有修复问题可以显式加入
\begin{equation}
 f(x)-b+1\le0.
 \label{eq:strict-repair}
\end{equation}
由于目标为整数，式~\eqref{eq:strict-repair} 与 $f(x)<b$ 等价。修复副本保留全部硬约束，将改善目标转化为一次可行性搜索。该局部搜索的失败不能证明 $Q_b$ 不可满足。

\subsection{由约束依赖扩展单坐标动作}
设一次探测希望修改坐标集 $U$。投影过程保护这些请求坐标，尝试通过受影响的约束恢复其他坐标。对于含一次可解变量的等式，仅当代数修复具有整数解且位于搜索表示范围内时执行更新；对于可识别的常量条件定义，则根据守卫真值同步更新输出。投影只执行有限轮，不能修复的部分交由后续搜索承担。

例如，在 $x+y=4$ 下最小化 $x^2+y^2$。从 $(0,4)$ 单独改变一个坐标会违反等式，而联合动作 $(0,4)\to(1,3)$ 保持可行，并将目标从 $16$ 降为 $10$。这一例子说明，动作的评价需要考虑其可恢复的联合变化。

候选生成还包括联合非零因子修改、目标前沿探测和差分约束联动。差分修复将目标探测作为固定要求，寻找其余坐标的线性见证；在包含析取的情形下，若当前文字选择导致负环，则尝试替代的文字见证。这些过程均只产生候选，完整约束核验决定候选能否更新 $(a,b)$。

\subsection{因子覆盖下的线性化修复}
代数投影难以直接处理一般非线性关系。我们进一步选择一组固定因子，使其余变量可以在严格改善条件下共同调整。设硬约束和目标中的单项式为 $m(x)=\prod_i x_i^{e_i}$。固定坐标集 $Z$ 满足
\begin{equation}
 \forall m,\qquad \sum_{i\notin Z}e_i\le1.
 \label{eq:cover}
\end{equation}
给定固定坐标上的整数试点 $v$，得到限制问题
\begin{equation}
 F[x_Z:=v]\land f[x_Z:=v]<b.
 \label{eq:restricted-lia}
\end{equation}
式~\eqref{eq:cover} 同时覆盖硬约束和目标的全部单项式，因此式~\eqref{eq:restricted-lia} 属于线性整数算术（LIA）。对于二分双线性因子图，可通过颜色类构造固定集；更一般的多项式则选择满足次数条件的覆盖。试点包括当前赋值附近的值、小域组合、零值和较远临界点。

\begin{proposition}[覆盖线性化]
\label{prop:cover}
若 $Z$ 对硬约束和目标中的每个单项式均满足式~\eqref{eq:cover}，则任意整数 $v$ 下，式~\eqref{eq:restricted-lia} 的算术原子与目标均为仿射表达式。其满足赋值补上 $x_Z=v$ 后满足原 $F$，且目标小于 $b$。
\end{proposition}
\begin{proof}
固定变量的整数幂仅贡献常数系数，每项剩余总次数至多为一，故代入后的多项式均为仿射式。代入保留原 CNF 的布尔连接和全部原子，补回固定坐标即可恢复原问题的满足赋值及其严格改善目标。
\end{proof}

这一变换得到原问题的精确限制，而非过近似。线性局部搜索器在限制问题上允许其余坐标联合修复；返回赋值补齐固定坐标后，仍需在完整非线性问题上重新检查。有限试点失败或后端未找到解，均不构成最优性证明。这里的线性局部搜索用于寻找单个改善赋值，后文的方向局部搜索则寻找能够永久保持可行的整数方向，两者承担不同角色。

\begin{algorithm}[htbp]
\caption{\textsc{CoupledImprove}：联合改善搜索的一段工作}
\label{alg:coupled}
\begin{algorithmic}[1]
\Require $F,f,a,b$，探索状态 $\sigma$，软优先级 $H$，本片预算 $\tau$
\State 恢复 $\sigma$，独立保留可行参照 $a$
\While{本片预算未耗尽}
 \State 若 $\sigma$ 可行则优先选择目标相关坐标，否则选择冲突子句相关坐标
 \State 根据 $H$ 调整候选坐标的扫描顺序
 \State 生成临界点和目标探测，并对关联定义执行有限轮投影
 \State 按适用条件加入联合因子、差分联动和前沿候选
 \State 按可行性、修复代价和目标排序，执行选定动作
 \If{得到完整可行且严格改善的候选 $c$} \State \Return $\mathsf{Candidate}(c)$ \EndIf
 \State 若达到修复时机，则续搜严格界副本或覆盖线性化问题
 \If{修复得到完整可行改善候选 $c$} \State \Return $\mathsf{Candidate}(c)$ \EndIf
 \State 若停滞，则尝试中性移动、目标扰动或重启修复
\EndWhile
\State 保留 $\sigma$ 及可续传的修复状态；\Return $\mathsf{NoConclusion}$
\end{algorithmic}
\end{algorithm}

算法~\ref{alg:coupled} 将上述动作统一到一段可恢复搜索中。中性移动保持目标与可行性，并降低固定的约束压力；扰动则允许暂时离开可行域，但不会覆盖独立保存的 $a$。一次改善被接纳后，还可沿未标记为定义变量的整数坐标联合变化做倍增外推，每个点重新投影和检查。有限次外推是候选生成手段，不是永久可行方向的证明。

\FloatBarrier
\section{冲突核心驱动的邻域扩展与反馈}
\label{sec:core}

小邻域能够降低一次改善查询的难度，却可能将所有更优解排除在外。我们利用可撤销的坐标锚点描述当前邻域，并用不可满足核心识别阻碍改善的一组固定条件。核心既指导邻域扩展，也为后续局部搜索提供变量优先级。

\subsection{锚定当前可行见证}
对坐标 $i$ 建立锚点 $\kappa_i:(x_i=a_i)$。设仍被锚定的坐标集合为 $A$，检查
\begin{equation}
 \Psi_{b,A}=F\land f<b\land\bigwedge_{i\in A}\kappa_i.
 \label{eq:anchored}
\end{equation}
完整 $F$ 和严格目标界是查询的固定部分，锚点通过 assumptions 传入。初始化释放定义性计算坐标及一个目标相关坐标，其余坐标尽量保持当前值。锚点可以作用于整数或布尔坐标，原 CNF 始终完整保留。

若式~\eqref{eq:anchored} 可满足，则得到一个严格改善的邻域候选；候选仍需通过完整精确检查。若不可满足，则从 assumptions 中提取锚点核心 $K\subseteq A$，使得
\begin{equation}
 F\land f<b\land\bigwedge_{i\in K}\kappa_i\quad\text{不可满足}.
 \label{eq:anchor-core}
\end{equation}

\begin{proposition}[核心指导的释放]
\label{prop:core}
若式~\eqref{eq:anchor-core} 成立，则任意满足 $F\land f<b$ 的赋值 $c$ 都至少改变一个核心坐标：存在 $i\in K$，使 $c_i\ne a_i$。
\end{proposition}
\begin{proof}
若 $c_i=a_i$ 对全部 $i\in K$ 成立，则 $c$ 同时满足式~\eqref{eq:anchor-core} 的全部合取项，与不可满足性矛盾。
\end{proof}

这一结论不要求核心最小，也不表示某个指定坐标必须改变。算法从非空 $K$ 中选择一个坐标释放，在后续时间片检查扩大的邻域。释放只是在已有逻辑信息支持下扩大搜索空间，并不保证紧接着就能找到改善。

\subsection{跨时间片扩展与软信息回灌}
算法~\ref{alg:core} 保存当前界下的锚点集和查询状态。一段工作最多执行一次实际 SMT 查询；查询构建或结果处理尚未完成时，可以在下一个时间片继续。新目标值会使旧邻域失效。

\begin{algorithm}[htbp]
\caption{\textsc{CoreNeighborhood}：核心指导的邻域演化}
\label{alg:core}
\begin{algorithmic}[1]
\Require $F,f,a,b$，邻域状态 $\eta$，本片预算 $\tau$
\State 若 $\eta$ 不属于当前界，则建立 $Q_b$ 并初始化锚点集 $A$
\If{$\eta$ 已结束} \State \Return $\mathsf{NoConclusion}$ \EndIf
\State 在本片预算内续做查询构建或尚未完成的结果处理
\State 若已可查询，则至多执行一次 $\Call{Check}{Q_b,\{\kappa_i:i\in A\}}$
\If{回答 SAT}
 \State 提取并精确检查完整赋值 $c$
 \If{$c\models F$ 且 $f(c)<b$} \State \Return $\mathsf{Candidate}(c)$ \EndIf
\ElsIf{回答 UNSAT}
 \State 提取锚点核心 $K$，将其变量作为软优先级返回
 \State 若 $K\ne\varnothing$，则释放其中一个坐标；否则结束当前邻域
\ElsIf{实际查询回答 UNKNOWN}
 \State 结束当前界的邻域尝试
\EndIf
\State 若本片仅在构建或结果处理时耗尽，则保留 $\eta$
\State \Return 已获得的软优先级与 $\mathsf{NoConclusion}$
\end{algorithmic}
\end{algorithm}

将核心涉及的变量加入集合 $H$，还能够改变后续 SLS 的候选扫描顺序。同一份冲突信息由此具有两种用途：在 SMT 一侧撤销部分锚点，在局部搜索一侧提高相关变量的探测优先级。$H$ 只是软提示，不是附加硬约束；接纳新目标值后会清空，避免将旧界冲突解释错误地沿用到新界。

当前实现遇到实际 UNKNOWN 时结束该界的邻域尝试，因为没有新的核心支持进一步释放。空核心虽然可表明搜索表示的严格改善问题不可满足，但这一接口仅结束邻域，不直接发布全局最优；有限最优结论统一通过原式证明通道确认。该机制以增强搜索为目的，但依赖 SMT 核心反馈，不属于完全独立于 SMT 的纯 SLS。

\FloatBarrier
\section{可行见证驱动的混合认证}
\label{sec:certificate}

有限最优与无界需要不同的证据。前者要求排除全部严格改善赋值，后者要求给出可以无限改善的一族可行赋值。我们首先说明原式严格改善查询如何与局部搜索交换模型，再介绍从可行见证合成整数射线的认证方法。

\subsection{原式严格改善检查与模型交换}
有限最优性通过 $Q_b^s=F_s\land f_s<b$ 检查。若查询 UNSAT，则不存在目标小于 $b$ 的原式赋值；结合当前可行目标值和式~\eqref{eq:values}，得到有限最优。若查询 SAT，则从原式模型提取目标并验证严格改善，再为搜索表示补齐辅助变量。补全成功且完整精确检查通过后，模型才可更新 $(a,b)$。

模型补全可能因坐标表示范围或预算限制失败，但不会改变原式 SAT 的逻辑结论。因此，本文采用的可续传证明配置在当前界已被 SAT 反驳后，暂停同界 UNSAT 尝试，直到出现新的可行目标值。SAT 的记录独立于候选接纳，避免因模型回灌失败而反复证明同一个已经被反驳的命题。

原式检查可以结合适用的守卫多项式放松、凸外近似、有限覆盖或原式 SMT 策略。使用这些过程时，必须区分结论的转移方向：必要条件放松的 UNSAT 在变换有效时可排除原式改善；放松的 SAT 则通常不能作为原式模型。类似地，局部固定下的 UNSAT 不能排除全局改善。

证明进度绑定当前目标界，相同 $b$ 下可以暂停和恢复；新界出现时重建相关状态。未取得结论时可以延长后续证明片，搜索再次改善后则缩短证明片，并在证明之间留出搜索窗口。连续证明窗口内接纳 SAT 改善后，可以立即检查新界，但只能使用该窗口剩余时间。调度本身不改变查询语义，也不会延长单例总预算。

\subsection{以当前模型选择子句见证}
有限多个改善点不能证明无界。我们从一个可行赋值 $a$ 出发，寻找对所有非负整数参数均可行、且目标永久改善的射线。

从每个子句 $C_j$ 中选取一个在 $a$ 下为真的文字 $\ell_j$，去重后得到集合 $W$，满足
\begin{equation}
 a\models\bigwedge_{\ell\in W}\ell,
 \qquad \bigwedge_{\ell\in W}\ell\ \Longrightarrow\ F.
 \label{eq:witness}
\end{equation}
选择时偏向较易处理的布尔或线性文字，但必须覆盖所有子句。随后选择固定因子集 $Z$，保持 $x_Z=a_Z$，使 $W$ 中的算术多项式和目标在代入后均为仿射式。这是精确代入后的仿射化，不能用某点的线性近似代替。

算法保留一组静态固定规则，并利用搜索中的改善活动、实际移动和违约压力优先保留活跃坐标自由。方向搜索失败后，将违约行映射回原始多项式因子，调整固定代价并交换锚点；新的固定集仍须通过精确仿射性检查。这些反馈只改变候选选择，不成为证明前提。某个固定方式失败，不表示原问题有界。与第~\ref{sec:coupled}~节的修复不同，这里只需维持一个覆盖全部子句的文字见证，但必须证明沿整条射线这些文字始终为真。

\subsection{整数方向约束与证书核验}
引入整数方向 $d$，定义
\begin{equation}
 x(t)=a+td,\qquad t\in\mathbb Z_{\ge0}.
 \label{eq:ray}
\end{equation}
式~\eqref{eq:ray} 只作用于整数坐标；布尔真值沿射线保持不变。固定整数坐标须满足 $d_i=0$（$i\in Z$）。对每个精确代入后为仿射的算术多项式 $p$，有
\begin{equation}
 p(a+td)=p(a)+t\,s_p(d),
 \label{eq:slope}
\end{equation}
其中 $s_p(d)$ 是关于整数方向的线性表达式。表~\ref{tab:slopes} 给出在基点为真的文字保持永久为真的充分条件。对于负等式，方向始终留在基点所在的正或负分支，不跨越零。

\begin{table}[htbp]
\centering
\caption{基点为真的见证文字及其永久保持条件。}
\label{tab:slopes}
\begin{tabular}{ll}
\toprule
见证文字或基点分支 & 方向约束 \\
\midrule
$p=0$ & $s_p(d)=0$ \\
$p\le0$ & $s_p(d)\le0$ \\
$\neg(p\le0)$，即 $p\ge1$ & $s_p(d)\ge0$ \\
$p\ne0$，且 $p(a)\ge1$ & $s_p(d)\ge0$ \\
$p\ne0$，且 $p(a)\le-1$ & $s_p(d)\le0$ \\
\bottomrule
\end{tabular}
\end{table}

最小化还要求目标斜率满足
\begin{equation}
 s_f(d)\le-1.
 \label{eq:objective-slope}
\end{equation}
以上条件组成 QF\_LIA 问题，由直接处理任意精度整数行的局部搜索合成方向。搜索先以齐次等式的整数联合移动保持变量间的耦合关系，再对其余约束进行加权违约修复和重启。找到满足方向约束的赋值后，使用精确整数运算重新核对基点可行性、子句覆盖、代入后的仿射性、固定坐标、各文字斜率及目标严格改善。未找到方向只表示本次搜索没有成功，不提供 UNSAT 结论。

\begin{algorithm}[htbp]
\caption{\textsc{WitnessRay}：从可行见证合成无界证书}
\label{alg:ray}
\begin{algorithmic}[1]
\Require $F,f,a,b$，本片预算 $\tau$
\State $W\gets$ 每个子句中选取一个在 $a$ 下为真的文字，并去重
\For{每个可在剩余预算内尝试的固定集 $Z$}
 \State 将 $x_Z$ 固定为 $a_Z$
 \If{$W$ 或 $f$ 的残余多项式仍非仿射} \State \textbf{continue} \EndIf
 \State $L\gets$ 固定整数坐标方向为零的约束；全部布尔真值保持不变
 \State $L\gets L\land W$ 的极性斜率条件 $\land\,(s_f(d)\le-1)$
 \State $(\mathit{ans},d,\rho)\gets\Call{DirectionSLS}{L,\text{本片预算}}$
 \If{$\mathit{ans}=\mathsf{SAT}$}
  \State $R\gets(a,d,Z,W)$
  \If{$\Call{ExactRayCheck}{F,f,R}$} \State \Return $\mathsf{VerifiedRay}(R)$ \EndIf
 \EndIf
 \State 若未成功，则根据违约反馈 $\rho$ 调整后续固定集；剩余预算可用于续搜
\EndFor
\State \Return $\mathsf{NoConclusion}$
\end{algorithmic}
\end{algorithm}

\begin{proposition}[整数射线证书的可靠性]
\label{prop:ray}
设 $a\models F$、$b=f(a)$，$W$ 满足式~\eqref{eq:witness}，固定 $x_Z=a_Z$ 后所选算术多项式及目标均为仿射。若整数方向 $d$ 在全部固定坐标上为零，全部布尔真值不变，且满足表~\ref{tab:slopes} 及式~\eqref{eq:objective-slope}，则对每个整数 $t\ge0$，$a+td\models F$，并且 $f(a+td)\to-\infty$。
\end{proposition}
\begin{proof}
固定整数坐标和布尔真值沿射线不变，因此精确代入得到的仿射式始终适用。由基点文字为真及式~\eqref{eq:slope}，表~\ref{tab:slopes} 的斜率条件保证每个所选算术文字对所有 $t\ge0$ 保持为真；特别地，负等式不会跨越零。式~\eqref{eq:witness} 保证每个子句始终至少有一个真文字，从而整式可行。目标满足 $f(a+td)=b+t\,s_f(d)$，由 $s_f(d)\le-1$ 得到其趋于负无穷。
\end{proof}

式~\eqref{eq:values} 将这一无界的目标值集合对应到原式；它不要求为每个 $t$ 显式恢复全部原式变量。作为说明，考虑 $x\ge0,y\ge0,z=xy$，最小化 $-z$。任意正整数 $k$ 给出基点 $(k,1,k)$ 和方向 $(0,1,k)$，于是
\[
 x(t)=k,\quad y(t)=1+t,\quad z(t)=k(1+t),\quad f(x(t))=-k(1+t).
\]
固定 $x=k$ 后，关系与目标沿方向均为仿射。这里证明的是全部参数值上的可行性与永久改善，而不是以一个有限大数替代无界。

\subsection{适用范围与结论边界}
有限大目标只用于触发和辅助探索。当前实现达到有利方向的阈值后，至多启动一次有界预算的射线构造；无界结论来自命题~\ref{prop:ray} 的核验条件，而非目标值的绝对大小。证书覆盖在选定文字见证和部分因子固定后可仿射化的整数射线，并不覆盖所有无界 NIA 问题，构造失败也不表示有界。

整体算法区分三种证据：完整精确检查提供可行目标值，原式严格改善问题不可满足提供有限最优，经过核验的射线提供无界性。搜索失败、局部邻域 UNSAT、固定集试探失败和预算耗尽均不会升级为全局认证。可靠性依赖前述导入语义前提和核验条件，算法不声称对一般 OMT(NIA) 完备。


\end{document}
